\section{Unique channel transitions and exact small-shell minimizers}
\label{sec:new-results}

Corollary~\ref{cor:spectrum-integrated} expresses the directional coefficients through
three explicit integrals. We now determine the ordering of these coefficients for every positive
channel parameter and, when the normal coefficient is strictly least,
identify the exact optimizer on every sufficiently small Procrustes shell.
The radii obtained here are uniform in the ambient dimension, but are not
evaluated numerically.

\subsection{A tilted Gaussian family}

For $m\ge3$ and $\mu\in\R$, set
\begin{align}
 w_{m,\mu}(c)&=e^{\mu c}\phi(c)^2\Phi(c)^{m-2},&
 Z_m(\mu)&=\int_\R w_{m,\mu}(c)\,dc,\label{new:tilt}\\
 a_m(\mu)&=Z_m(\mu)^{-1}\int_\R h(c)^2w_{m,\mu}(c)\,dc,&
 b_m(\mu)&=Z_m(\mu)^{-1}\int_\R \sigma(c)w_{m,\mu}(c)\,dc,\notag\\
 H_m(\mu)&=Z_m(\mu)^{-1}\int_\R h(c)w_{m,\mu}(c)\,dc.\notag
\end{align}
Here $q=\phi/\Phi$, $h(c)=c+q(c)$ and $\sigma=1-qh$, as in
Theorem~\ref{thm:spectrum}. The normalizer is positive and finite.
At $\mu=\lambda/\sqrt\alpha$,
\[
 A_\lambda=C_\lambda a_m(\mu),\qquad
 B_\lambda=C_\lambda b_m(\mu),\qquad
 C_\lambda=\frac{\mu e^{-\mu^2/2}}m Z_m(\mu)>0.
\]
Consequently the normalized decoding curvatures are
\begin{equation}\label{new:normalized-curvatures}
 L_0=\frac{m-1}{m}a_m(\mu),\qquad
 L_1=\frac{a_m(\mu)+(m-2)b_m(\mu)}2,\qquad
 L_\perp=\frac{m-1}{2}.
\end{equation}
The endpoint $\mu=0$ below is an endpoint of the normalized moment family;
it is not division by $C_\lambda$ at a zero-amplitude channel.

\begin{theorem}[Unique ordered channel crossings]\label{thm:ordered-switches}
Let $m\ge4$. There are unique positive parameters
\[
 0<\mu_{1\perp}(m)<\mu_{0\perp}(m)<\mu_{01}(m)
\]
satisfying, respectively,
\[
 a_m+(m-2)b_m=m-1,\qquad a_m=m/2,\qquad a_m=mb_m.
\]
Each equation has no other root in $[0,\infty)$. For $\lambda>0$,
with $\mu=\lambda/\sqrt\alpha$, the four open intervals have the orderings
\begin{align*}
 0<\mu<\mu_{1\perp}&:\quad \Lambda_0<\Lambda_1<\Lambda_\perp,\\
 \mu_{1\perp}<\mu<\mu_{0\perp}&:\quad \Lambda_0<\Lambda_\perp<\Lambda_1,\\
 \mu_{0\perp}<\mu<\mu_{01}&:\quad \Lambda_\perp<\Lambda_0<\Lambda_1,\\
 \mu>\mu_{01}&:\quad \Lambda_\perp<\Lambda_1<\Lambda_0.
\end{align*}
At the three endpoints the indicated pair agrees, and the remaining
inequality is strict. In particular
$\min\{\Lambda_0,\Lambda_\perp\}<\Lambda_1$ for every $\lambda>0$.
The thresholds in $\lambda$ are $\sqrt\alpha$ times these thresholds.
\end{theorem}

\begin{proof}
Appendix~\ref{ap:transitions} proves continuity and strict increase of
$a_m$, $b_m$, $H_m$ and $a_m/b_m$, as well as
\[
 \frac12<b_m(\mu)<1\quad(m\ge4,\ \mu\ge0),\qquad
 2a_m+(m-1)b_m=m+\mu H_m.
\]
At zero the identity gives $a_m<m/2$ and, using $b_m<1$,
$a_m+(m-2)b_m<m-1$. Using $b_m>1/2$ also gives $a_m<mb_m$.
Moreover $2a_m(\mu)>1+\mu H_m(0)$, so $a_m$ tends to infinity.
The three continuous strictly increasing functions
$a_m+(m-2)b_m$, $a_m$, and $a_m/b_m$ therefore cross their targets once.
At $a_m=m/2$, the bound $b_m>1/2$ puts the first crossing to the left
and the third to the right. Finally
$L_1-L_0=(m-2)(mb_m-a_m)/(2m)$ gives the displayed sign chart.
The analytic ingredients, including the strict benchmark at $m=4$,
are proved in full in Appendix~\ref{ap:transitions}.
\end{proof}

\begin{theorem}[The triangular channel crossing]\label{thm:triangular-switch}
For $m=3$ there is a unique $\mu_*(3)>0$ such that
$\mu_*H_3(\mu_*)=1$, equivalently $a_3(\mu_*)+b_3(\mu_*)=2$.
For $0<\mu<\mu_*$ one has $\Lambda_1<\Lambda_\perp$; for
$\mu>\mu_*$ the reverse inequality holds. They agree only at $\mu_*$.
The space $\mathcal H_0$ is absent.
\end{theorem}
\begin{proof}
The balance identity gives
$(a_3+b_3)/2=(3+\mu H_3)/4$.
On $[0,\infty)$, the function $\mu H_3(\mu)$ is continuous and strictly
increasing from zero to infinity, since $H_3$ is positive and increasing.
This proves the assertion. Its integral ingredients are included in
Appendix~\ref{ap:transitions}.
\end{proof}

\subsection{The exact strict-normal shell optimum}

For $m\ge3$ and $\lambda>0$, impose the strict-normal condition
\begin{equation}\label{new:normal-gap}
 \begin{cases}
  \mu>\mu_{0\perp}(m),&m\ge4,\\
  \mu>\mu_*(3),&m=3,
 \end{cases}
 \qquad \mu=\lambda/\sqrt\alpha.
\end{equation}
Equivalently $a_m(\mu)>m/2$ for $m\ge4$, and
$a_3(\mu)+b_3(\mu)>2$ for $m=3$. These are precisely the conditions that
the normal coefficient be strictly smaller than every in-plane coefficient
whose space is present.

\begin{theorem}[Exact optimizer on small Procrustes shells]\label{thm:normal-shell}
Assume \eqref{new:normal-gap}. There is $\rho_0=\rho_0(m,\lambda)>0$
such that, for every $d\ge m$, every unit regular simplex
$V=[v_1\ \cdots\ v_m]$ in $\R^d$, and every $0<\rho<\rho_0$, the minimum
of $\delta_\lambda(X)=\Pc(G_\Delta;\lambda)-\Pc(G_X;\lambda)$ over all
unit configurations with $d_P(X,V)=\rho$ is attained precisely by the
orthogonal images of
\begin{equation}\label{new:normal-minimizer}
 N_i=\kappa v_i+\sqrt{1-\kappa^2}\,z,\qquad
 \kappa=1-\frac{\rho^2}{2m},\qquad
 \|z\|=1,\quad z\perp\operatorname{span}\{v_i\}.
\end{equation}
The normal $z$ may be chosen once for all these radii. Different choices
of $z$ give the same orthogonal orbit. The minimal deficit is exactly
\begin{equation}\label{new:exact-deficit}
 \delta_{\min}(\rho)=
 \Pc(G_\Delta;\lambda)-
 \Pc\left(G_\Delta;\lambda\left(1-\frac{\rho^2}{2m}\right)\right).
\end{equation}
\end{theorem}

\begin{proof}
Appendix~\ref{ap:normal} supplies the complete geometric argument.
The centered square root
$T=(\alpha PG_XP)^{1/2}$ satisfies, on a uniform neighborhood,
$d_P(X,V)^2=2m-2\tr(T)/\alpha$.
On the shell $\tau=\rho^2$, write $T=\kappa P+D$, where
$D$ is symmetric, centered and trace-free.
The covariance reconstructed from $(\tau,D)$ defines a smooth
extension of the decoding deficit. Permutation symmetry makes $D=0$
stationary for every small $\tau$.
Its Hessian at $(0,0)$ is positive definite under \eqref{new:normal-gap};
the precise coefficients and the square-root coordinate normalization
are derived in Appendix~\ref{ap:normal}.
Continuity then gives a strict quadratic lower bound near this ray.
Every sufficiently small actual shell lies in that neighborhood.
Thus $D=0$ is its unique minimizing covariance.
It is realized by \eqref{new:normal-minimizer}, with exact distance $\rho$,
and equality of Gram matrices is equivalent to an orthogonal change
of realization. A common Gaussian shift in the normal direction gives
\eqref{new:exact-deficit}. The neighborhood estimates are made in the
fixed $m\times m$ Gram space, so $\rho_0$ is independent of $d$ and $V$.
\end{proof}

\begin{theorem}[Quartic expansion with a sixth-order remainder]
\label{thm:normal-taylor}
Let $p(\lambda)=\Pc(G_\Delta;\lambda)$. This function is real analytic.
For fixed $m\ge3$ and $\lambda>0$ satisfying \eqref{new:normal-gap},
there are $C=C(m,\lambda)>0$ and $r=r(m,\lambda)>0$ such that
Theorem~\ref{thm:normal-shell} holds for $0<\rho<r$ and
\begin{align}
 \delta_{\min}(\rho)
 &=\frac{\lambda p'(\lambda)}{2m}\rho^2
   -\frac{\lambda^2p''(\lambda)}{8m^2}\rho^4+\mathcal R_{m,\lambda}(\rho),
 \label{new:quartic}\\
 |\mathcal R_{m,\lambda}(\rho)|&\le C|\rho|^6
       \qquad (|\rho|<r).\label{new:remainder}
\end{align}
For negative $\rho$, the deficit in this display denotes the even
extension given by the right side of \eqref{new:exact-deficit}.
All its odd derivatives at zero vanish. In particular there is no cubic
term. The leading coefficient satisfies
\begin{equation}\label{new:curvature-check}
 \frac{\lambda p'(\lambda)}{2m}=\Lambda_\perp
       =\frac{m-1}{2}C_\lambda.
\end{equation}
\end{theorem}

\begin{proof}
All real exponential moments of the Gaussian maximum in \eqref{eq:pc}
are finite, so its moment generating function, and hence $p$, is real
analytic. Taylor's theorem through the quadratic term at $\lambda$
has error $O(|s-\lambda|^3)$.
Substituting $s=\lambda-\lambda\rho^2/(2m)$ gives
\eqref{new:quartic}--\eqref{new:remainder}; the sign of the quartic term
comes from subtracting $p(s)$ from $p(\lambda)$.
The exact expression is even in $\rho$, which gives the assertion about
all odd derivatives. The normal covariance path has initial velocity
$(J-I_m)/m$ in the squared radius, so its first deficit derivative is
$C_\lambda(m-1)/2$. Differentiating the exact value gives the other side
of \eqref{new:curvature-check}. Detailed domination and remainder
arguments are in Appendix~\ref{ap:normal}.
\end{proof}

The equal-curvature boundary is excluded, and no claim of a finite-radius
in-plane optimizer is made. The positive radius and remainder constant
are existential; they do not supply an effective numerical certificate.
The channel crossings, the strict-normal optimization theorem and the
quantified Taylor expansion are formalized over the actual Gaussian
integrals and orthogonal Procrustes distance. Appendix~\ref{ap:lean}
records the exact scope, theorem correspondence and reproducible Lean
checks; the other global bounds and numerical experiments in this paper
are not included in that formalization.
